\section{创新点、资料整合与演进路线}
\label{sec:synthesis}
\subsection{作品的创新点}
嘟嘟哒的核心设计主线是\textbf{自适应}：以当前讨论与任务为输入，在管理员设定的范围内选择参与方式、业务插件和计算路由，并通过 WebUI 呈现实际结果。自适应分别落实到能力选择、计算选择和交互选择，形成可解释的群聊服务机制。

\textbf{自适应插件选择。}系统综合讨论主题、多人参与、需求持续性与当前请求，在管理员预先授权的只读服务中选择并启用匹配插件。以选课讨论为例，群成员逐步提出讲解方式、证明训练和课业负担等条件，当前请求达到触发条件后，系统启用评课插件并进入课程资料查询。能力选择因此与连续讨论关联，减少用户记忆命令和手动切换服务的负担；启用原因与实际能力调用可在运行记录中核对。实现机制见第 5.4、7 节。

\textbf{任务驱动的自适应路由。}感知结果提供意图与复杂度信号，档位策略结合置信度、可用模型和预算选择本轮计算配置，路由器再校验角色、能力与输出容量。简单交流、校园查询和复杂比较可以使用不同档位；推理强度与可见篇幅分别控制，避免仅按答案长短分配计算资源。自适应体现在运行选择随任务改变，选择依据可记录并在 WebUI 中查看。当前策略由规则和配置定义，其工程目标是匹配任务所需的计算投入；尚未完成同条件成本与效果对照实验。机制见第 6 节。

\textbf{自适应群聊参与与表达。}系统根据提及、追问、持续讨论和群策略决定回复、查询、澄清或静默，再结合任务要求组织回答篇幅与人格表达。参与选择与冷却、频率限制分别执行，既利用多人交流中的情境信号，也保持管理员对参与节奏的控制。工具事实与 Persona 分层，使严谨解释和日常表达能够在同一运行路径中按需组织。机制见第 5、8 节。

\textbf{WebUI 工作台设计。}工作台将真实 QQ 会话与 Agent 对话、运行、配置置于统一界面，并在桌面并排布局与手机单页切换中保留当前账号和会话。管理员按群设置偏好、允许范围与锁定项，运行后查看实际模型、插件、上下文覆盖、耗时和选择原因，再调整策略。QQ 与 Web 预览共享核心，构成“配置—预览—观察—调整”的连续操作，使自适应机制既能被配置，也能被检查。详细设计见第\ref{sec:webui}节。

这些创新属于校园群聊场景下的系统设计与交互机制创新。当前自适应由语义感知、确定性规则、管理员配置与运行状态共同实现；在线 Bandit、长期记忆学习和策略更新另有独立研究与接入计划。

\subsection{本地资料汇总范围}
本次扫描主仓、主仓阶段记录、两套历史原型/插件目录、原提交材料、Web 试用报告、嘟嘟哒文档站、相关部署记录、录制基准，以及培养方案数据说明。共登记 441 份当前主仓及独立来源文档；加上工作树与发布副本为 8,944 份物理文件，按逻辑路径归并为 446 个条目。数值反映文档整理范围，不是独立研究或测试数量。

完整清单见同目录的“本地文档汇总.md”，每项记录标题、主要议题、采用方式与副本数量；\code{source-inventory.json} 保留各副本位置。报告主体按主题综合材料，避免照搬按时间累积的开发日志。私有运行配置、聊天数据库、凭据、第三方依赖文档和批量培养方案数据条目不作为技术报告输入。

\begin{longtable}{@{}L{29mm}L{62mm}L{58mm}@{}}
\caption{文档体系的整合方式}\\\toprule
资料类型 & 主要贡献 & 在本报告中的使用\\\midrule\endfirsthead
\toprule 资料类型 & 主要贡献 & 在本报告中的使用\\\midrule\endhead
设计文档与 ADR & 核心依赖、消息与状态、模型、能力、记忆、人格、安全、部署边界 & 第 3—\ref{sec:deployment}节；以当前组合根验证已接入内容。\\
研究报告 & 语义感知、群体情境、主动参与、记忆检索、离线策略与插件化方向 & 第 5—9、\ref{sec:synthesis}节；保留为设计动机或后续方向，不转成运行能力。\\
重构与阶段记录 & 从早期插件组合到统一智能体核心的迁移、接口收敛、真实纵向接入 & 解释历史路径、类名与职责变化；旧的“当前状态”按日期使用。\\
服务与插件 README & 服务输入输出、认证和缓存、依赖、显式命令与默认开关 & 第 4、7、\ref{sec:webui}—\ref{sec:deployment}节；区别独立插件、可选 MCP 与 Planner 能力。\\
评测、录制与试用 & 固定消息、真实查询、回答审校、模拟入站、UI 问题及修复 & 第\ref{sec:evaluation}节；分别保留执行模式、日期、失败项和样本限制。\\
文档站与历史原型 & 早期模块设计、13 阶段流水线、8 角色路由、开发与部署经验 & 对照采用接口分工思想；不移植旧原型的功能和测试数字。\\
原提交 PDF / Word & 作品定位、演示配套信息、旧报告与简介 & 统一成熟技术语体，补充架构、API、验证和完成度说明。\\\bottomrule
\end{longtable}

\subsection{资料中的主要差异及处理}
\textbf{旧状态与新实现。}文档站的架构页面固定在早期提交，\code{current-state.md} 也包含重构前基线；其中“尚未实现”或旧插件清单不能覆盖当前组合根。新版按现有源码描述运行路径，历史状态用于解释演进。

\textbf{原型与主仓。}历史 \code{dududa20-prototype} 和 \code{dududa20-plugin} 记录过多步规划、八角色路由、个人信息及记忆等功能，并有自己的测试数字。当前主仓采用不同的实现、组合和阶段划分，新报告不把它们视为同一发布物。

\textbf{设计目标与组件接入。}完整 Memory、主动日报、多 OC 和 Skill 组合在多份资料中反复出现。源码存在、离线验证、配置登记、正式接入与长期效果是不同证据，当前状态已在第\ref{sec:evaluation}节统一。

\textbf{模型配置与测试口径。}Luna/Terra/Sol 是档位别名，旧材料中还使用过其他模型；本次演示记录使用 DeepSeek。常规推理强度与录制时关闭 thinking 的配置分别注明，不把测试数量或样本时延外推到全部模型。

\subsection{下一阶段的技术路线}
近期优先改进来源覆盖、查询前澄清、近期摘要的时间与归属，分解等待时间，并对现有广覆盖题集完成同范围的真实模型复测。参与机制需要在获准的自然群窗口上评估时机、回应对象、误触发和重复建议，而不只评估是否成功生成文字。

随后接入 Memory 的显式写入、检索、删除与反馈闭环，在真实但经授权的数据上比较无记忆、近期优先和 BM25。多 OC 与群友印象应在有明确资料来源和修改方式的基础上建设。

较长远的扩展是以现有业务能力为基础的多步 Skill，以及同角色、同任务范围内的路由优化。二者均应从具体任务和可测收益出发，保持当前业务的可维护性，并保留用户与管理员对群聊参与的控制。
\source{[S1—S18]；文档全集与副本索引随 LaTeX 源码交付。}
